\cchapter{Synchronization}{synchronization}
\label{chap:synchronization}

The \kcode{barrier} construct is a stand-alone directive that requires all threads
of a team (within a contention group) to execute the barrier and complete
execution of all tasks within the region, before continuing past the barrier.

The \kcode{critical} construct is a directive that contains a structured block. 
The construct allows only a single thread at a time to execute the region.
Multiple \kcode{critical} regions may exist in a parallel region, and they may act
cooperatively (only one thread at a time in all \kcode{critical} regions) or
separately (only one thread at a time in each \kcode{critical} regions when a
unique name is supplied on each \kcode{critical} construct). An optional (lock)
\kcode{hint} clause may be specified on a named \kcode{critical} construct to
provide the OpenMP runtime guidance in selecting a locking mechanism.

On a finer-grained scale, the \kcode{atomic} construct allows only a single
thread at a time to have atomic access to a storage location involving a single
atomic read, write, or update statement, and a limited number of combinations
when specifying the \kcode{capture} and \kcode{compare} \plc{extended-atomic}
clauses.  Unlike the \kcode{read} and \kcode{write} \plc{atomic} clauses, the
\kcode{update} \plc{atomic} clause is optional and implied for update
statements if not explicitly specified.. The \plc{memory-order} clause can be
used to specify the degree of memory ordering enforced by an \kcode{atomic}
construct. From weakest to strongest, they are \emph{relaxed} (the default),
\emph{acquire} and/or \emph{release} (specified with \kcode{acquire},
\kcode{release}, or \kcode{acq_rel} clauses), and \emph{sequentially consistent}
(specified with \kcode{seq_cst}).  Please see the details in the \docref{atomic
Construct} subsection of the \docref{Synchronization Constructs and Clauses}
chapter in the OpenMP Specification document.

The \kcode{ordered} construct may be block-associated or stand-alone. The
block-associated form may appear in worksharing-loop (\kcode{for} or
\kcode{do}), \kcode{simd}, or worksharing-loop SIMD (\kcode{for simd} or
\kcode{do simd}) region. It sequentializes and orders the execution of the
\kcode{ordered} regions while allowing code outside the region to run in
parallel. The stand-alone \kcode{ordered} construct specifies cross-iteration
dependences in a \emph{doacross} loop nest. The \kcode{doacross} clause uses a
\kcode{sink} \plc{dependence-type}, along with an iteration vector argument
(\plc{vec}) to indicate the iteration that satisfies the dependence.  The
\kcode{doacross} clause with a \kcode{source} \plc{dependence-type} (and
optional \kcode{omp_curr_iteration} keyword as the iteration vector
representing the current iteration) specifies dependence satisfaction.

The \kcode{flush} directive is a stand-alone construct for enforcing consistency
between a thread's view of memory and the view of memory for other threads (see
the Memory Model chapter of this document for more details). When the construct
is used with an explicit variable list, a \emph{strong flush} that forces a
thread's temporary view of memory to be consistent with the actual memory is
applied to all listed variables. When the construct is used without an explicit
variable list and either without a \plc{memory-order} clause or with the
\kcode{seq_cst} \plc{memory-order} clause, a strong flush is applied to all
locally thread-visible data as defined by the base language, and additionally
the construct provides both acquire and release memory ordering semantics.
When an explicit variable list is not present and a \plc{memory-order} clause
other than \kcode{seq_cst} is present, the construct provides acquire and/or
release memory ordering semantics according to the \plc{memory-order} clause,
but no strong flush is performed. A resulting strong flush that applies to a
set of variables effectively ensures that no memory (load or store) operation
for the affected variables may be reordered across the \kcode{flush} directive.

General-purpose routines provide mutual exclusion semantics through locks, 
represented by lock variables.  
The semantics allows a task to \emph{set}, and hence 
\emph{own} a lock, until it is \emph{unset} by the task that set it. A 
\emph{nestable} lock can be set multiple times by a task, and is used
when in code requires nested control of locks.  A \emph{simple lock} can
only be set once by the owning task. There are specific calls for the two
types of locks, and the variable of a specific lock type cannot be used by the
other lock type.  

Any explicit task will observe the synchronization prescribed in a
\kcode{barrier} construct and an implied barrier.  Also, additional
synchronizations are available for tasks.  A task will wait at a
\kcode{taskwait} for all of its child tasks to complete.  A \kcode{taskgroup}
construct creates a region in which the current task is suspended at the end of
the region until all child tasks, and their descendants, have completed.
Scheduling constraints on task execution can be prescribed by the
\kcode{depend} clause to enforce a dependence on previously generated tasks.
More details on controlling task executions can be found in the \docref{Tasking
Constructs} chapter in the OpenMP Specifications document. %(DO REF. RIGHT.)


%===== Examples Sections =====
\input{synchronization/critical}
\input{synchronization/worksharing_critical}
\input{synchronization/barrier_regions}
\input{synchronization/atomic}
%\pagebreak
\section{Atomic Compare}
\label{sec:cas}

\index{constructs!atomic@\kcode{atomic}}
\index{atomic construct@\kcode{atomic} construct}
\index{clauses!capture@\kcode{capture}}
\index{clauses!compare@\kcode{compare}}
\index{capture clause@\kcode{capture} clause}
\index{compare clause@\kcode{compare} clause}

The \kcode{compare} clause was added to the \plc{extended-atomic} clauses
for C/C++ in OpenMP 5.1 and extended to Fortran in OpenMP 6.0.
The \kcode{compare} clause extends the semantics to perform the \kcode{atomic}
update conditionally. 

In the following example, two formats of structured blocks
are shown for associated \kcode{atomic} constructs with the \kcode{compare} clause.
In the first \kcode{atomic} construct, the format forms a \emph{conditional update} statement.
In the second \kcode{atomic} construct the format forms a \emph{conditional expression} statement.
The ``greater than'' and ``less than'' forms and 
the conditional expression form from Fortran 2023
are available to the \kcode{compare} clause for Fortran in OpenMP 6.0.
% One can use the \vcode{max} and \vcode{min} functions with the \kcode{atomic update}
% construct to perform the C/C++ example operations.

\cexample[5.1]{cas}{1}
\ffreeexample[6.0]{cas}{1}
\clearpage

In OpenMP 5.1 the \kcode{compare} clause with the \kcode{capture} clause
was added to support \emph{Compare And Swap} (CAS) semantics 
in C/C++ and extended to Fortran in OpenMP 6.0.
In the following example, the \ucode{enqueue} routine
(a naive implementation of a Michael and Scott enqueue function), uses the
\kcode{compare} with \kcode{capture} clause to perform compare
(\plc{x == e}) and swap (\plc{if-else} assignments) of the
form in C/C++: 
\begin{description}[noitemsep,labelindent=5mm,widest=f90]
\item \splc{{ r = x == e; if (r) { x = d; } else { v = x; } }}.
\end{description}

The equivalent form for Fortran pointers is:
\begin{description}[noitemsep,labelindent=5mm,widest=f90]
\item \splc{r = associated(x, e); if (r) then; x => d; else; v => x; end if}
\end{description}
where the intrinsic function \bcode{associated(\plc{x, e})} is used
for pointer comparison.
In the example code, the corresponding variables for ``\plc{x, e, d, v}''
are C/C++ variables ``\ucode{queue->tail}, \ucode{queue->next}, 
\ucode{node}, \ucode{node->next}'', 
or Fortran variables ``\ucode{queue\%tail}, \ucode{queue\%next}, 
\ucode{node}, \ucode{node\%next}'', respectively.
The example program concurrently enqueues nodes from an array of nodes
(\ucode{nodes[N]} or \ucode{nodes(N)}).

\cexample[5.1]{cas}{2}
\ffreeexample[6.0]{cas}{2}

\input{synchronization/atomic_restrict}
\input{synchronization/atomic_hint}
%\pagebreak
\section{Synchronization Based on Acquire/Release Semantics}
\label{sec:acquire_and_release_semantics}

\index{flushes!acquire}
\index{flushes!release}
\index{clauses!memory ordering clauses}
\index{memory ordering clauses!acquire@\kcode{acquire}}
\index{acquire clause@\kcode{acquire} clause}
\index{memory ordering clauses!release@\kcode{release}}
\index{release clause@\kcode{release} clause}
\index{memory ordering clauses!acq_rel@\kcode{acq_rel}}
\index{acq_rel clause@\kcode{acq_rel} clause}
\index{flush construct@\kcode{flush} construct}
\index{atomic construct@\kcode{atomic} construct}
\index{clauses!acquire@\kcode{acquire}}
\index{clauses!release@\kcode{release}}
\index{clauses!acq_rel@\kcode{acq_rel}}

As explained in the Memory Model chapter of this document, a flush operation
may be an \emph{acquire flush} and/or a \emph{release flush}, and OpenMP 5.0
defines acquire/release semantics in terms of these fundamental flush
operations.  For any synchronization between two threads that is specified by
OpenMP, a release flush logically occurs at the source of the synchronization
and an acquire flush logically occurs at the sink of the synchronization.
OpenMP 5.0 added memory ordering clauses -- \kcode{acquire}, \kcode{release}, and
\kcode{acq_rel} -- to the \kcode{flush} and \kcode{atomic} constructs for
explicitly requesting acquire/release semantics.  Furthermore, implicit flushes
for all OpenMP constructs and runtime routines that synchronize OpenMP threads
in some manner were redefined in terms of synchronizing release and acquire
flushes to avoid the requirement of strong memory fences (see the \docref{Flush
Synchronization and Happens Before} and \docref{Implicit Flushes} sections of the
OpenMP Specification document).

The examples that follow in this section illustrate how acquire and release
flushes may be employed, implicitly or explicitly, for synchronizing threads. A
\kcode{flush} directive without a list and without any memory ordering clause
can also function as both an acquire and release flush for facilitating thread
synchronization.  Flushes implied on entry to, or exit from, an atomic
operation (specified by an \kcode{atomic} construct) may function as an acquire
flush or a release flush if a memory ordering clause appears on the construct.
On entry to and exit from a \kcode{critical} construct there is now an implicit
acquire flush and release flush, respectively.

%%%%%%%%%%%%%%%%%%

\index{constructs!critical@\kcode{critical}}
\index{critical construct@\kcode{critical} construct}
\index{flushes!implicit}
The first example illustrates how the release and acquire flushes implied by a
\kcode{critical} region guarantee a value written by the first thread is visible
to a read of the value on the second thread. Thread 0 writes to \ucode{x} and
then executes a \kcode{critical} region in which it writes to \ucode{y}; the write
to \ucode{x} happens before the execution of the \kcode{critical} region,
consistent with the program order of the thread.  Meanwhile, thread 1 executes a
\kcode{critical} region in a loop until it reads a non-zero value from
\ucode{y} in the \kcode{critical} region, after which it prints the value of
\ucode{x}; again, the execution of the \kcode{critical} regions happen before the
read from \ucode{x} based on the program order of the thread. The \kcode{critical}
regions executed by the two threads execute in a serial manner, with a
pair-wise synchronization from the exit of one \kcode{critical} region to the
entry to the next \kcode{critical} region.  These pair-wise synchronizations
result from the implicit release flushes that occur on exit from
\kcode{critical} regions and the implicit acquire flushes that occur on entry to
\kcode{critical} regions; hence, the execution of each \kcode{critical} region in
the sequence happens before the execution of the next \kcode{critical} region.
A ``happens before'' order is therefore established between the assignment to \ucode{x}
by thread 0 and the read from \ucode{x} by thread 1, and so thread 1 must see that
\ucode{x} equals 10.

%\pagebreak
\cexample[5.0]{acquire_release}{1}
\ffreeexample[5.0]{acquire_release}{1}

\index{constructs!atomic@\kcode{atomic}}
\index{atomic construct@\kcode{atomic} construct}
\index{atomic construct@\kcode{atomic} construct!write clause@\kcode{write} clause}
\index{atomic construct@\kcode{atomic} construct!read clause@\kcode{read} clause}
\index{atomic construct@\kcode{atomic} construct!memory ordering clauses}
\index{write clause@\kcode{write} clause}
\index{read clause@\kcode{read} clause}
\index{clauses!write@\kcode{write}}
\index{clauses!read@\kcode{read}}
\index{memory ordering clauses!seq_cst@\kcode{seq_cst}}
\index{seq_cst clause@\kcode{seq_cst} clause}
\index{clauses!seq_cst@\kcode{seq_cst}}
In the second example, the \kcode{critical} constructs are exchanged with
\kcode{atomic} constructs that have \emph{explicit} memory ordering specified. When the
\emph{atomic read} operation on thread 1 reads a non-zero value from \ucode{y}, this
results in a release/acquire synchronization that in turn implies that the
assignment to \ucode{x} on thread 0 happens before the read of \ucode{x} on thread
1. Therefore, thread 1 will print ``x = 10''.

\cexample[5.0]{acquire_release}{2}
\ffreeexample[5.0]{acquire_release}{2}

%\pagebreak
\index{constructs!atomic@\kcode{atomic}}
\index{atomic construct@\kcode{atomic} construct!relaxed atomic operations}
\index{flush construct@\kcode{flush} construct}
In the third example, \kcode{atomic} constructs that specify relaxed atomic
operations are used with explicit \kcode{flush} directives to enforce memory
ordering between the two threads. The explicit \kcode{flush} directive on thread
0 must specify a release flush and the explicit \kcode{flush} directive on
thread 1 must specify an acquire flush to establish a release/acquire
synchronization between the two threads. The \kcode{flush} and \kcode{atomic}
constructs encountered by thread 0 can be replaced by the \kcode{atomic} construct used in
Example 2 for thread 0, and similarly the \kcode{flush} and \kcode{atomic}
constructs encountered by thread 1 can be replaced by the \kcode{atomic}
construct used in Example 2 for thread 1.

\cexample[5.0]{acquire_release}{3}
\ffreeexample[5.0]{acquire_release}{3}

Example 4 will fail to order the write to \ucode{x} on thread 0 before the read
from \ucode{x} on thread 1. Importantly, the implicit release flush on exit from
the \kcode{critical} region will not synchronize with the acquire flush that
occurs on the \emph{atomic read} operation performed by thread 1. This is because
implicit release flushes that occur on a given construct may only synchronize
with implicit acquire flushes on a compatible construct (and vice-versa) that
internally makes use of the same synchronization variable. For a
\kcode{critical} construct, this might correspond to a \emph{lock} object that
is used by a given implementation (for the synchronization semantics of other
constructs due to implicit release and acquire flushes, refer to the \docref{Implicit
Flushes} section of the OpenMP Specifications document).  Either an explicit \kcode{flush}
directive that provides a release flush (i.e., a flush without a list that does
not have the \kcode{acquire} clause) must be specified between the
\kcode{critical} construct and the \emph{atomic write}, or an atomic operation that
modifies \ucode{y} and provides release semantics must be specified.

\cexample[5.0]{acquire_release_broke}{4}
\ffreeexample[5.0]{acquire_release_broke}{4}

\input{synchronization/ordered}
\input{synchronization/depobj}
\input{synchronization/doacross}
\input{synchronization/locks}
\input{synchronization/init_lock}
\input{synchronization/init_lock_with_hint} 
\input{synchronization/lock_owner}
\input{synchronization/simple_lock}
\input{synchronization/nestable_lock}
%\pagebreak
\section{\kcode{safesync} Clause}
\label{sec:safesync}

\index{clauses!safesync@\kcode{safesync}}
\index{safesync clause@\kcode{safesync} clause}

Parallel regions in OpenMP are usually executed by a team of threads that
should be able to freely take \emph{divergent code paths} and synchronize with
each other using various point-to-point synchronization mechanisms. However, on
certain target architectures in which OpenMP threads may instead execute a
single stream of (possibly predicated) instructions in lock step, two threads
may not be able to synchronize with each other from logically divergent code
without a resulting deadlock. The \kcode{safesync(\plc{width})} clause may be
specified on a \kcode{parallel} construct to control which threads in a team
should be able to synchronize with each other from divergent code paths. The
effect is that the team is partitioned into ``progress groups'' of consecutive
threads (in terms of thread number) of size \plc{width}, except for the final
group which can have less than \plc{width} threads. Threads in the same
progress group may execute a single sequence of instructions, while threads in
different progress groups will execute distinct sequences of instructions. Thus,
two threads in different progress groups can safely synchronize from divergent
code paths.

If the \kcode{safesync} clause isn't present the behavior is as if the
\plc{width} is 1 for a \kcode{parallel} construct encountered on the host
device or when the \kcode{device_safesync} requirement has been specified for
that compilation unit. In other words, any two threads in the team can
synchronize from divergent code paths by default under those conditions.

The following example uses a ticket lock implementation in a \kcode{target}
region that requires thread synchronization from divergent code paths. Under
the progress guarantees for threads that are defined in the 6.0 specification,
this code can deadlock without an explicit use of the \kcode{safesync} clause
when executing on a non-host device. For example, on an NVIDIA GPU, the
implementation may map threads in the team to individual GPU threads in a warp
which does not support such synchronization between divergent threads. To make
this work, the \kcode{safesync} clause (with the omitted \plc{width}
defaulting to 1) can be specified. The example includes two variants of a
\kcode{parallel} directive, one with a \kcode{safesync} clause and one without,
via a metadirective that is conditioned on whether the device is a host device
or non-host device. Alternatively, the \kcode{device_safesync} requirement
could be specified with a \kcode{requires} directive, and this would require
\kcode{safesync} behavior as a default even on a non-host device.

\cexample[6.0]{safesync}{1}
\ffreeexample[6.0]{safesync}{1}



